%%%%

%%%   Самарский государственный технический университет

%%%   Кафедра прикладной математики и информатики

%%%

%%%   Преамбула для статьи в журнал "Вестник Самарского государственного

%%%   технического университета. Серия: «Физико-математические науки»"

%%%   Ориентирована под MikTEx 2.4 for Win32

%%%

%%%   Хотя сам журнал собирается под GNU/Linux на дистрибутиве TeXLive



\documentclass[11pt,a4paper,twoside]{article}



\usepackage[cp1251]{inputenc}

\usepackage[T2A]{fontenc}

\usepackage[english,russian]{babel}

\usepackage{samgtubib}

\usepackage[dvips]{graphicx}

\usepackage[multiple]{footmisc}



\usepackage{samgtu}

\renewcommand{\VestnikYear}{2009} % Год издания Вестника

\renewcommand{\VestnikNumber}{2\,(19)} % Номер Вестника

\renewcommand{\VestnikMonth}{Ноябрь} % Месяц выхода Вестника

\renewcommand{\VestnikData}{01.11.09} % Дата подписания Вестника в печать

\renewcommand{\VestnikUPL}{26,64} % Усл. печ. л.

\renewcommand{\VestnikUIL}{25,73} % Уч.-изд. л.

\renewcommand{\VestnikRegNumber}{479/09} % Регистрационный номер

\renewcommand{\VestnikOrder}{994} % Номер заказа







%\usepackage{nccboxes}

%%

%\begin{document}

%%

%%





\renewcommand{\firstpage}{188}

\renewcommand{\lastpage}{190}





\renewcommand{\baselinestretch}{0.8}



%\ShortPartition{Химическая физика}{}



%\udk{534.222.2} \msc{??????}

\udk{532.539:536.711} \msc{80A30}





\title{О методе оценки латуни как материала для облицовки боеприпасов типа <<ударное ядро>> }

{О методе оценки латуни как материала для облицовки боеприпасов  \dots}



\author{И.\,И.~Реут, А.\,Л.~Кривченко}{Р\,е\,у\,т~И.\,И., К\,р\,и\,в\,ч\,е\,н\,к\,о~А.\,Л.}



\institute{\samgtu.}



\email{E-mails}{enterfax@mail.ru, snex@rambler.ru}



\otherinf{\fullauthor{Игорь Игоревич Реут}  \position{аспирант} \dept{каф. защиты в чрезвычайных ситуациях}

\fullauthor{Александр Львович Кривченко} \regalia{д.т.н., доц.} \position{профессор} \dept{каф. защиты в чрезвычайных ситуациях}}



\keywords{ударная адиабата, латунь, экспресс-метод, торцевое метание}



\workabstract{Проведён расчёт скорости торцевого метания пластин зарядами A-IX-1 и ТГ~50/50 по динамическим адиабатам метаемых элементов. 

Показано, что кинетический модуль для единицы массы при метании пластины из латуни зарядом A-IX-1  примерно на 4\,\% выше и на 8\,\% выше при её метании зарядом ТГ~50/50, чем при метании медной пластины.

}



\engtitle{On assessment method of brass as the material for explosively formed penetrator lining}



\engauthor{I.\,I.~Reout, A.\,L.~Krivchenko}{R\,e\,o\,u\,t~I.\,I., K\,r\,i\,v\,c\,h\,e\,n\,k\,o~A.\,L.}





\enginstitute{\engsamgtu.}



\otherenginf{\fullauthor{Igor I. Reout}

\position{Postgraduate Student} \dept{Dept. of Protection of Extremal Situation} \fullauthor{Alexander~L.~Krivchenko} \regalia{Dr. Sci. (Techn.)} \position{Professor} \dept{Dept. of Protection of Extremal Situation}}





\engabstract{The calculation of a front tossing speed of plates by A-IX-1 and Cyclotol-50/50 charges was made on the basis of tossing elements dynamic adiabats. It is shown that as opposed to copper plate the kinetic module for the mass unit is about 4\,\% higher for tossing the brass plate in case of using A-IX-1 and is about 8\,\% higher in case of using Cyclotol-50/50.}



\engkeywords{shock adiabat, brass, express method, front tossing}



\date{27/IX/2011}

\revisiondate{13/X/2011}





\maketitle



\small



\Section{Введение} В настоящее время одной из важнейших задач для разработки боеприпасов является поражение бронированных целей. 

В~этом плане параметры поражающих факторов с~бронёй (преградой) определяются ударной адиабатой соударяющихся тел.

Одним из вопросов разработки систем кинетического оружия является разработка ударника и~определение его свойств. 

Это также актуально для зарядов взрывного бурения, ударных возбудителей сейсмосигнала и~кумулятивных перфораторов. 

В~настоящее время наиболее изученным в~качестве ударников являются медь (ударные ядра), сталь и вольфрам (бронебойные и~подкалиберные снаряды). 

Характеристики же композиционных сплавов в~качестве ударников изучены неполно, хотя многие, особенно на основе d-металлов,  являются наиболее перспективными в этом плане.



Сложные динамические взаимодействия представляют самостоятельный интерес для различных областей науки и~техники, в~том числе и~для разработки новых систем динамического оружия. В~настоящее время не существует теоретических методов оценки эффективности действия композиционных материалов для облицовки боеприпасов типа <<ударное ядро>>. Подбор материалов ограничен экспериментальными данными в основном по меди и~стали. Поэтому в настоящей работе рассматривается латунь в качестве облицовки боеприпасов типа <<ударное ядро>>.



Критерием оценки эффективности действия металла может стать кинетический модуль $E$ для единицы массы облицовки:

\[

E=W^2/2,

\]

где $W$ "--- скорость метания облицовки.

Для этого необходимо разработать с достаточно высокой достоверностью теоретический экспресс"=метод оценки скорости

метания, основанный на ударно"=волновом анализе.



\smallskip



\Section{Определение скорости торцевого метания пластины по динамическим адиабатам метаемых элементов} 

Наиболее удобная форма задания ударной адиабаты в  координатах <<$D$~-- $U$>> имеет вид

\begin{equation}

D=\alpha +\lambda U,

\label{reut2:eq1}

\end{equation}

где $D$ "--- скорость ударной волны, $U$ "--- массовая скорость, $\alpha $ и $\lambda $ "--- эмпирические коэффициенты. 

К~настоящему времени константы $\alpha $ и $\lambda $ определены для многих материалов. 

Константа $\alpha $ численно равна объёмной скорости звука $C$, соответствующей начальной адиабатической объёмной сжимаемости вещества.









Величину $\lambda =dD/dU$ в \eqref{reut2:eq1} можно рассматривать как производную  в~некоторой промежуточной точке ударной адиабаты. 

Значение $\lambda $ меняется вдоль адиабаты Гюгонио с~увеличением плотности вещества в ударной волне \cite{reut2:bib1}. В работе \cite{reut2:bib2} на ос\-нове анализа экспериментальных ударных адиабат, взятых в основном из

\cite{reut2:bib1}, показано, что ударные адиабаты элементов составляют две группы с

одинаково меняющимися со степенью сжатия вещества во фронте ударной волны

величинами~$\lambda $.





\begin{wrapfigure}[47]{O}{.5\textwidth}

\centering





\includegraphics[width=.38\textwidth]{./00PIC/p-reut1}

\caption{Схема определения параметров ударной волны, входящей  в

металлическую пластину, и скорости её метания в первом импульсе: \scriptsize

точка 1 "--- параметры детонации на адиабате продуктов детонации;

точка 2 "--- зеркальное отражение адиабаты продуктов детонации на данную

адиабату металлической пластины;

$2U$ "--- скорость метания пластин при выходе ударной волны на свободную

поверхность}



\small 



\vspace{2mm}



\setlength{\extrarowheight}{1pt}

\setlength{\tabcolsep}{2.5pt}

{\bf \footnotesize Характеристики латуни и её компонентов}

\begin{tabular}{c|c|c|c|c} \hline

\footnotesize Металл & \footnotesize  $\rho$, г/см$^{3}$ & \footnotesize $\lambda $ & \footnotesize $C$, м/с & \cbox{\footnotesize Массовая \\ \footnotesize доля } \\ \hline

\rule{0mm}{12pt}%

Zn&7,14&1,45&3300&0,235 \\ 

Cu&8,96&1,48&3980&0,765 \\ 

Cu+Zn&8,63&1,51&3725& "--- \\ \hline

\end{tabular}





\vspace{2mm}



\includegraphics[width=.48\textwidth]{./00PIC/reut2.eps}

\caption{Ударные адиабаты:  \scriptsize 1 "--- экспериментальная ударная адиабата меди

\cite{reut2:bib5}; 2 "--- экспериментальная ударная адиабата цинка \cite{reut2:bib5};

3 "--- экспериментальная адиабата латуни \cite{reut2:bib1}; 4 "--- расчётная адиабата латуни

}

\end{wrapfigure}





Принято считать, что соотношение (\ref{reut2:eq1}) может служить уравнением состояния

химических элементов и сложных веществ, включая и сплавы \cite{reut2:bib3}. К тому же

линейные соотношения в координатах <<$D$~-- $U$>>, описывающие ударную адиабату в области более высоких

давлений, будут иметь произвольные коэффициенты $\alpha $ и $\lambda $.







Для расчёта ударных адиабат при ударно"=волновом взаимодействии необходимым

условием является точный прогноз значения объёмной скорости звука в~композиционном материале, в~который входит ударная волна. Согласно данным \cite{reut2:bib4},

 объёмная скорость звука, даже если она не определена экспериментально,

с~высокой степенью точности рассчитывается по правилу симметричного 8"=элементного

окружения данного элемента в~таблице периодического закона Д.\,И.~Менделеева. 

Согласно данному правилу объёмная скорость звука $C$ определяется как

среднеарифметическое всех восьми окружающих элементов:

\[

C=\frac{1}{8} \sum \limits_{i=1}^8 C_i,

\]

где $C_i$ "--- объёмная скорость звука элементов окружения.



Аналогичным методом определяется коэффициент $\lambda$:

\[

\lambda =\frac{1}{8} \sum \limits_{i=1}^8  \lambda_i,

\]

где $\lambda_i$ "--- значения коэффициентов у элементов окружения.



Расчёт массовой скорости $U$ для ударной адиабаты осуществляется по методу акустического приближения и скорости метания в первом отражении ударных волн. Схематически метод представлен на рис.~1.







Данные для расчёта скорости звука и коэффициента ударной адиабаты $\lambda $ на примере латуни приведены в таблице.





Расчётная и экспериментальная ударные адиабаты для латуни и её компонентов в

координатах <<$D$~-- $U$>>  приведены на рис.~2, из которого видно, что имеет место близость хода экспериментальной и вычисленной ударных адиабат на основе расчётных значений объёмной скорости звука и коэффициента ударной адиабаты  $\lambda $, полученных методом <<правила симметричного окружения>>, и практически полного их совпадения в диапазоне массовых скоростей $U$ в от 700 до 1600~м/с.

Именно в этом диапазоне лежат массовые скорости материала пластины, 

которые для меди и латуни при их нагружении зарядом A-IX-1 составляют 955~м/с и 987~м/с соответственно, 

а при нагружении этих облицовок зарядом ТГ~50/50 "--- 730 м/с и 763~м/с соответственно. 

Кинетический модуль для единицы массы при метании пластины из латуни зарядом А-IX-1  примерно на 4\,\% выше и на 8\,\% выше при её метании зарядом ТГ~50/50, чем при метании медной пластины.



\smallskip

\Section{Выводы}

Разработанный экспресс"=метод оценки эффективности действия \linebreak взрывчатых

веществ вполне применим для металлических материалов поражающих элементов

боеприпасов типа <<ударное ядро>> и даёт возможность определить наиболее

оптимальный материал, хотя и требует экспериментальной проверки.



\begin{thebibliography}{9}



\Bibitem{reut2:bib1}

\by van Thiel~M., Kusubov~A.\,S., Mitchell~A.\,C.

\preprint Compendium of shock wave data

\preprintinfo  Technical Report UCRL-50108

\publaddr Livermore, CA

\publ Lawrence Radiation Lab.

\yr 1967



\RBibitem{reut2:bib2}

\by Анисичкин~В.\,Ф.

\paper Обобщённые ударные адиабаты элементов

\jour ПМТФ

\yr 1978

\issue 3

\pages 117--121

\transl англ. пер.:~

\jour J. Appl. Mech. Tech. Phys.

\vol 19

\issue 3

\pages 376--379

%DOI: 10.1007/BF00850824

\paper Generalized shock adiabats of the elements

\by Anisichkin~V.\,F.





\RBibitem{reut2:bib3}

\by Андреев~С.\,Г.,  Бабкин~А.\,В., Баум~Ф.\,А. и др.

\book Физика взрыва

\bookinfo в 2-x томах

\ed Л.~П.~Орленко

\bookvol 1

\publaddr М.

\publ Физматлит

\yr 2004

\totalpages 823

\misctransl

\by Andreev~S.\,G., Babkin~A.\,V., Baum~F.\,A., et al.

\book Physics of Explosion

\ed L.~P.~Orlenko 

\bookvol 1

\publ Fizmatlit

\publaddr Moscow

\yr 2004

\totalpages 823



\RBibitem{reut2:bib4}

\paper Расчет ударных адиабат d-металлов и их сплавов с использованием

периодического закона Д.\,И.~Менделеева

\by Кривченко~А.\,Л., Кривченко~Д.\,А., Реут~И.\,И., Чуркин~О.\,Ю.

\inbook Ударные волны в конденсированных средах

\publaddr СПб.

\yr 2008

\pages 253--255

\misctransl 

\by Krivchenko~A.\,L., Krivchenko~D.\,A., Reout~I.\,I.,  Churkin~O.\,Yu.

\paper Calculation of shock adiabats for d-metals and their alloys on the basis of Mendeleev periodic law

\inbook Shock waves in condensed matter

\publaddr St. Petersburg

\yr 2008

\pages 253--255



\RBibitem{reut2:bib5}

\by Иванченко~Е.\,С.

\thesis Прецизионные модели ударных адиабат и база ТЕФИС 

\thesisinfo Автореф. дисс.~\dots\ канд. физ.-мат. наук: 05.13.18

\publ ИММ РАН

\publaddr М.

\yr 2009

\totalpages 22

\misctransl

\by Ivanchenko~E.\,S.

\thesis Precise models of shock adiabatic and TEFIS database

\thesisinfo Ph. D. Thesis (Phys. \& Math.)

\publaddr Moscow

\publ IMM RAN

\yr 2009

\totalpages 22



\end{thebibliography}





\noindent

\hskip6.5cm \thedate \par\noindent

\hskip6.5cm\therevdate





\vspace{-10mm}







\makeengtitle



\renewcommand{\baselinestretch}{0.9}



 \newpage

%

% \tableofengcontents

%

%\end{document}

